\documentclass[11pt]{article}
\usepackage{amsmath,amssymb,amsthm}
\usepackage{mathtools}
\usepackage[margin=1.1in]{geometry}
\usepackage{hyperref}
\usepackage{cleveref}

\theoremstyle{plain}
\newtheorem{theorem}{Theorem}[section]
\newtheorem{proposition}[theorem]{Proposition}
\newtheorem{corollary}[theorem]{Corollary}
\theoremstyle{definition}
\newtheorem{definition}[theorem]{Definition}
\newtheorem{example}[theorem]{Example}

\DeclareMathOperator{\supp}{supp}
\DeclarePairedDelimiter{\abs}{\lvert}{\rvert}
\DeclarePairedDelimiter{\norm}{\lVert}{\rVert}
\newcommand{\N}{\mathbb{N}}
\newcommand{\R}{\mathbb{R}}

\title{Real Analysis II --- Lecture Notes 4\\\large Uniform Convergence}
\author{Notes by A.~Petrov}
\date{}

\begin{document}
\maketitle

\section{Pointwise and uniform convergence}

\begin{definition}\label{def:unif}
A sequence $(f_n)_{n\in\N}$ of functions $f_n\colon E\to\R$ converges \emph{uniformly}
to $f$ if for every $\varepsilon>0$ there is $N$ such that
\[
  \sup_{x\in E}\abs{f_n(x)-f(x)} < \varepsilon \quad\text{for all } n\ge N.
\]
\end{definition}

\begin{example}
Let $f_n(x) = x^n$ on $[0,1]$. Then $f_n\to f$ pointwise, where $f(x)=0$ for $x<1$ and
$f(1)=1$, but not uniformly, since $\norm{f_n-f}_\infty = 1$ for every $n$.
\end{example}

\begin{theorem}\label{thm:cont}
If $f_n\to f$ uniformly on $E$ and each $f_n$ is continuous at $x_0$, then $f$ is
continuous at $x_0$.
\end{theorem}

\begin{proof}
Let $\varepsilon>0$ and choose $N$ with $\norm{f_N-f}_\infty<\varepsilon/3$. Choose
$\delta>0$ such that $\abs{f_N(x)-f_N(x_0)}<\varepsilon/3$ whenever
$\abs{x-x_0}<\delta$. Then
\begin{align}
\abs{f(x)-f(x_0)}
 &\le \abs{f(x)-f_N(x)} + \abs{f_N(x)-f_N(x_0)} + \abs{f_N(x_0)-f(x_0)} \notag\\
 &< \tfrac{\varepsilon}{3}+\tfrac{\varepsilon}{3}+\tfrac{\varepsilon}{3} = \varepsilon. \label{eq:three}
\end{align}
\end{proof}

The argument in \eqref{eq:three} is the classical ``$\varepsilon/3$ trick''.

\begin{corollary}
The uniform limit of continuous functions is continuous.
\end{corollary}

\section{Integration and differentiation}

\begin{proposition}\label{prop:int}
If $f_n\to f$ uniformly on $[a,b]$ and each $f_n$ is Riemann integrable, then
\[
  \lim_{n\to\infty}\int_a^b f_n(x)\,dx = \int_a^b f(x)\,dx.
\]
\end{proposition}

By \cref{thm:cont} and \cref{prop:int}, uniform convergence preserves both
continuity and integrals. It does \emph{not} preserve derivatives:
\[
  f_n(x) = \frac{\sin(nx)}{\sqrt n} \xrightarrow{\text{unif.}} 0,
  \qquad\text{but}\qquad f_n'(0) = \sqrt{n}\to\infty.
\]

\begin{definition}
The \emph{support} of $f$ is $\supp f = \overline{\{x : f(x)\neq 0\}}$.
\end{definition}

\end{document}
